% AUTORIDAD DE FUENTE: Exposé IX, §§7.2.2--7.2.6,
% líneas 1940--1967; índices de escaneo 417--418, folios 406--407.
% Las capas francesas D/C no presentan diferencias materiales.

\medskip
\noindent
7.2.2. Apliquemos 7.2.1 al esquema formal en grupos $\widehat A$
considerado en 7.1, en el caso favorable tratado allí en el que el
esquema abeliano formal $\widehat B$ es algebraizable (cf. 7.1.5), y
supongamos además que el toro $T_o$ sobre $S_o$ es isotrivial. Obtenemos
así un esquema en grupos $A^\natural$ sobre $S$, llamado el
\underline{grupo de Raynaud} de $A$ sobre $S$, caracterizado salvo
isomorfismo por un isomorfismo canónico de esquemas formales en grupos
sobre $\operatorname{Specf}(V)=\widehat S$:

%%%%%%%%%%  scan idx 418  |  folio 407  %%%%%%%%%%

\begin{equation}\tag{7.2.3}
 (A^\natural)^\wedge\simeq\widehat A.
\end{equation}
Se deduce de 7.2.1 que
\begin{equation}\tag{7.2.4}
 \Phi=A^\natural/A^{\natural o}
\end{equation}
es un grupo finito étale sobre $S$, y que la componente neutra
\[
 A^{\natural o}=A^{o\natural}
\]
está provista de la estructura de una extensión
\begin{equation}\tag{7.2.5}
 0\longrightarrow T\longrightarrow A^{\natural o}
 \longrightarrow B\longrightarrow0,
\end{equation}
donde $B$ es el esquema abeliano sobre $S$ cuyo completado formal es el
esquema abeliano formal $\widehat B$ construido en 7.1, y $T$ está
caracterizado, salvo isomorfismo único, como el toro isotrivial sobre
$S$ que levanta $T_o$. Al restringir a los $S_n$, la estructura de
extensión (7.2.5) da las estructuras de extensión sobre
$A_n^{\natural o}=A_n^o$ definidas por (7.1.1).

\medskip
\noindent
7.2.6. También se deduce inmediatamente de 7.2.1 que el grupo de
Raynaud $A^\natural$, así como la sucesión exacta (7.2.5), dependen
funtorialmente de $A$ (que satisface las condiciones consideradas), y
que su formación conmuta con todo cambio de base $S'\longrightarrow S$
asociado a un homomorfismo de anillos topológicos ádicos
$V\longrightarrow V'$.
